\section{Findings \& Discussion of Results}
\label{sec:Fin}

This section reports the experimental outcomes of the server-side ECU trace replay with LARVA instrumentation toggled on/off. Table~\ref{tab:rv-overhead} uses the stock LARVA-generated monitor (no logging/stacktrace optimisations), with 1 warmup run and 5 measured runs per mode on the same machine and JVM conditions---matching the \emph{Original} variant summarised in Table~\ref{tab:rv-overhead-variants}. Additional discussion below references supplementary benchmark sessions with a tuned \texttt{loadMapMismatch} threshold (\texttt{|Load-MAPSource| > 60}), 3 warmup runs, and 10 measured runs per mode. The experiment assumes telemetry is transmitted from an IoT/ECU endpoint to a server-side verification API.

\subsection{Result presentation}

Two regimes were evaluated:
\begin{itemize}
    \item \textbf{Baseline regime:} clean trace (\texttt{20180713-home2mimos\_clean.csv}).
    \item \textbf{Abnormal regime:} injected trace (\texttt{20180713-home2mimos\_clean\_abnormal.csv}) with approximately 20\% violation rows across four flags.
\end{itemize}

Table~\ref{tab:rv-overhead} summarises core timing and throughput metrics. Standard deviation is reported with each average to capture run-to-run variability.
All elapsed-time values reported in this section correspond to full-trace processing per request (68{,}433 rows).

\begin{table}[ht]
    \centering
    \caption{Runtime verification overhead summary (unmodified generated monitor; 5 measured runs per mode)}
    \label{tab:rv-overhead}
    \scriptsize
    \resizebox{\columnwidth}{!}{%
    \begin{tabular}{l|c|c}
    \textbf{Metric} & \textbf{Baseline} & \textbf{Abnormal} \\ \hline
    RV ON elapsed (ms) & 125.8 $\pm$ 11.19 & 4824.4 $\pm$ 199.48 \\
    RV OFF elapsed (ms) & 118.2 $\pm$ 17.36 & 119.8 $\pm$ 13.85 \\
    RV ON throughput (rows/s) & 547{,}374.9 $\pm$ 47{,}727.4 & 14{,}204.4 $\pm$ 593.9 \\
    RV OFF throughput (rows/s) & 587{,}489.3 $\pm$ 72{,}679.2 & 577{,}078.3 $\pm$ 63{,}349.7 \\
    RV ON CPU (ms) & 128.1 $\pm$ 13.07 & 3562.5 $\pm$ 189.44 \\
    RV OFF CPU (ms) & 121.9 $\pm$ 27.95 & 125.0 $\pm$ 22.10 \\
    Overhead (\%) & 6.43 & 3927.05 \\
    \end{tabular}
    }
\end{table}

\subsection{Analysis and interpretation}

\textbf{RQ1 (detection of injected abnormalities).} The abnormal regime produced substantial RV ON flag totals (\texttt{coolantTooHigh}=17105, \texttt{voltageOutOfRange}=17100, \texttt{throttleSpike}=17110, \texttt{loadMapMismatch}=17115; 68{,}430 events in aggregate across the four flags), while RV OFF flags remained zero. This indicates that violation reporting is coupled to monitor execution as intended and that injected abnormal classes are detected consistently.

\textbf{RQ2 (monitoring overhead).} On the unmodified generated monitor, baseline overhead is moderate (6.43\%), while abnormal overhead is extremely large (3927.05\%), showing that violation-heavy execution paths dominate runtime when stock diagnostics (logging and stacktrace-backed state strings) are enabled. In other words, monitor cost is not constant; it is strongly state-dependent. These overhead values should be interpreted as server-side verification cost under a telemetry-upload architecture; Table~\ref{tab:rv-overhead-variants} shows how far overhead falls once those diagnostics are targeted.

\textbf{RQ3 (stability and threshold sensitivity).} Baseline runs show standard deviations that are modest relative to means; abnormal RV-on timings exhibit wider spread, consistent with allocation-heavy diagnostic paths. After tuning \texttt{loadMapMismatch} from 30 to 60, baseline false positives reduced (aggregate baseline \texttt{loadMapMismatch} decreased compared with pre-tuning runs), while abnormal detection remained strong. This supports threshold calibration as a practical control for false-positive pressure.

\subsection{Controlled root-cause experiment}

To isolate the source of abnormal-path overhead, an additional controlled comparison was executed on three monitor variants under identical benchmark settings (1 warmup, 5 measured runs): (i) original generated monitor, (ii) logging disabled, and (iii) logging disabled plus removal of bad-state stacktrace/string construction.

\begin{table}[ht]
    \centering
    \caption{Abnormal-path overhead across monitor variants (controlled 3-way comparison)}
    \label{tab:rv-overhead-variants}
    \scriptsize
    \resizebox{\columnwidth}{!}{%
    \begin{tabular}{l|c|c|c|c}
    \textbf{Variant} & \textbf{Base O/H (\%)} & \textbf{Abn O/H (\%)} & \textbf{Abn RV ON (ms)} & \textbf{Abn RV ON CPU (ms)} \\ \hline
    Original & 6.43 & 3927.05 & 4824.4 & 3562.5 \\
    Logging disabled & 14.61 & 683.79 & 851.2 & 850.0 \\
    No stacktrace & 6.75 & 9.39 & 125.8 & 125.0 \\
    \end{tabular}
    }
\end{table}

This progression shows that the dominant abnormal overhead originates from generated diagnostic behaviour in bad-state handling (especially stacktrace-backed state-string construction), not from core guard checks. CPU measurements followed the same trend as elapsed time, while heap-delta remained noisy due to garbage collection effects.

\subsection{Comparison with related work and quality considerations}

The observed trade-off aligns with prior RV literature: runtime monitoring provides actionable safety/cybersecurity visibility, but overhead depends on property design, trigger frequency, and deployment context~\cite{krichen2023survey,zhang2024context,jaksic2018algebraic}. The clear difference between baseline and abnormal runtime also mirrors broader automotive V\&V concerns, where worst-case and abuse-case behaviour must be measured rather than assumed from nominal conditions~\cite{ekert2021cybersecurity}.

\textbf{Validity.} Internal validity is strengthened by same-endpoint comparisons, warmup before measured runs, and on/off toggling as the only intended treatment factor. Construct validity is supported by explicit metrics (elapsed time, throughput, p95/p99, overhead) and deterministic replay inputs.

\textbf{Reliability.} Repeated measurements with reported standard deviations (5-run batches for Table~\ref{tab:rv-overhead}; 10-run batches where noted above) indicate reproducibility under stable local conditions.

\textbf{Generalizability.} External validity remains limited to this server-side replay implementation and selected property set. Results should not be extrapolated directly to production ECUs without hardware-in-the-loop or in-vehicle evaluation.